\chapter{Phase 0 - Staging of Project}

\section{Phase objective}

As the first stage of development, the repository’s base structure 
is established and the data contract is defined using the 
intermediate model, before implementing the system’s business 
logic. This design decision allows all stages of the pipeline 
to be developed using a previously specified and consistent 
communication interface.

Defining the contract early on ensures that components can evolve 
independently, provided they adhere to the agreed-upon data model.

\section{File structure}
The project structure is based on packages, each of which reflects 
one of the project phases and can be tested individually.
The following shows the organization of the repository in its 
main folder:

\begin{lstlisting}[language=bash, caption={Structure of \texttt{src/spekforge}}]
src/spekforge/
    models/    # shared intermediate model
    ingest/    # Phase 1 - Ingest and normalization
    extract/   # Phase 2 - LLM extraction
        rag/   # Module of RAG for extended documentation
    assemble/  # Phase 3 - Final OpenAPI assembly
    verify/    # Phase 4a - static and dynamic validation
    generate/  # Phase 4b - SDK,mock server, contract tests
    orchestrator/ # state proper machine
    evaluation/ # ground-truth metrics   
\end{lstlisting}

\section{Background configuration}
The project uses Python 3.12 with \texttt{uv}

\subsection{\texttt{pyproject.toml}}
\begin{lstlisting}[language={[x86masm]Assembler}, basicstyle=\ttfamily\footnotesize, caption={Project dependencies and config}]
[project]
name = "spekforge"
version = "0.1.0"
description = "Reconstruction of OpenAPI specifications from informal documentation, with loop verification"
readme = "README.md"
requires-python = ">=3.12"
dependencies = [
    "pydantic>=2.7",
    "pyyaml>=6.0",
    "markitdown>=0.0.1a3",
    "openapi-spec-validator>=0.7",
    "schemathesis>=3.30",
    "httpx>=0.27",
    "ollama>=0.3",
    "instructor>=1.4",
    "python-dotenv>=1.0",
    "typer>=0.12",
]

[project.optional-dependencies]
rag = [
    "chromadb>=0.5",
]

[dependency-groups]
dev = [
    "pytest>=8.0",
    "pytest-cov>=5.0",
    "ruff>=0.5",
]

[build-system]
requires = ["hatchling"]
build-backend = "hatchling.build"

[tool.pytest.ini_options]
testpaths = ["tests"]
\end{lstlisting}

A local execution solution was chosen as the language model provider 
in order to avoid the costs associated with using cloud services 
and to ensure the tool’s availability throughout the development 
and evaluation process. Specifically, \textbf{Ollama} was selected
as the model execution environment, complemented by Instructor, 
a library that allows for validating generated responses against 
Pydantic models and performing automatic retries when the output 
does not conform to the expected structure.

This design decision provides a fully local, reproducible experimentation 
environment that is independent of external services.

\subsection{\texttt{.gitignore}}
The following are the exclusions defined for the project's version 
control. These include files and directories generated during 
development, dependencies, temporary artifacts, and other elements 
that are not part of the source code and are therefore excluded 
from the repository according to the rules set forth 
in the .gitignore file.

\begin{lstlisting}[caption={version control exclusions}]
# Virtual Environment
.venv/
venv/

#uv
uv.lock.tmp

#Python
__pycache__/
*.py[cod]
*.egg-info/
.eggs/
build/
dist/

# testing and coverage
.pytest_cache/
.coverage
htmlcov/
.ruff_cache/
.mypy_cache/

# env files
.env

#outputs generated by pipeline
data/outputs/*
!data/outputs/.gitkeep

# Ground - truth data
tests/fixtures/*/ground_truth/*
!tests/fixtures/*/ground_truth/.gitkeep

# IDE
.vscode/
.idea/

#OS
.DS_Store
\end{lstlisting}

For the purpose of testing, the official OpenAPI specification for 
the target API is deliberately excluded from the inputs provided 
to the system. By keeping it hidden during pipeline execution, 
we ensure that the reconstruction is performed solely based on 
the available documentation and not through direct inspection of 
the official specification.

This decision ensures that the comparison performed at the end 
of the process constitutes an objective evaluation of the system’s 
ability to reconstruct the specification, avoiding biases resulting 
from prior access to the reference document.

\subsection{Dependencies installation}

To install the dependencies, the command \texttt{uv sync} was executed,
creating \texttt{uv.lock} and the virtual environment \texttt{.venv/}, 
so that the components could be tested properly. Figure 2.1 below
shows the resolution and installation of dependencies without 
errors

\realcapture{0.85}{Installation.png}{Command output of \texttt{uv sync} completing dependency resolution}

\section{Intermediate model}
To build the intermediate model, three Pydantic classes were defined,
which are distributed as \texttt{ParamSpec}, \texttt{SchemaSpec},
and \texttt{EndpointSpec}, in that order of dependency. 

\subsection{\texttt{ParamSpec}}

\begin{lstlisting}[language=Python, caption={\texttt{src/spekforge/models/param.py}}]
from enum import Enum
from pydantic import BaseModel

class ParamLocation(str, Enum):
    PATH = "path"
    QUERY = "query"
    HEADER = "header"
    COOKIE = "cookie"

class ParamSpec(BaseModel):
    name: str
    location: ParamLocation
    required: bool = False
    type: str = "string"
    description: str | None = None
    example: str | None = None
\end{lstlisting}
In this case, \texttt{ParamLocation} was defined as an \texttt{Enum}, 
rather than a free-form text string (\texttt{str}), in order to restrict 
the language model, during Phase 2, to selecting only values belonging to 
a previously defined and valid set. When generating the JSON schema 
using (\texttt{model\_json\_schema()}), this enumeration translates 
into an explicit restriction within the \textit{JSON Schema} used 
during the \textit{function calling}.

\subsection{\texttt{SchemaSpec}}

\begin{lstlisting}[language=Python, caption={\texttt{src/spekforge/models/schema.py}}]
from pydantic import BaseModel, Field

class SchemaSpec(BaseModel):
    name: str | None = None
    type: str = "object"
    format: str | None = None
    description: str | None = None
    properties: dict[str, "SchemaSpec"] = Field(default_factory=dict)
    required: list[str] = Field(default_factory=list)
    items: "SchemaSpec | None" = None
    enum: list[str] | None = None
    ref: str | None = None

SchemaSpec.model_rebuild()
\end{lstlisting}

The \texttt{SchemaSpec} model has a recursive structure, the \texttt{properties} attribute corresponds to a set of objects of the same \texttt{SchemaSpec} type, while the \texttt{items} field represents the schema of the elements contained in an array. This definition allows for the modeling of nested data structures, such as a \texttt{Pet} object that contains a list of \texttt{Tag} objects, while preserving the hierarchy of the OpenAPI specification.

Due to this self-reference, it is necessary to call the \texttt{model\_rebuild()} method once the class definition 
is complete. This procedure resolves references to the class 
itself, since, during the evaluation of the \texttt{SchemaSpec} 
body, the corresponding symbol has not yet been fully defined.

The \texttt{ref} field is intended for Phase 3; \texttt{assemble/dedup.py}
will use it to indicate that a schema has already been registered in the components
under the specified name. This is done to prevent the LLM from describing the same
object in slightly different ways in different sections of the documentation.

\subsection{\texttt{EndpointSpec}}

\begin{lstlisting}[language=Python, caption={\texttt{src/spekforge/models/endpoint.py}}]
from enum import Enum
from pydantic import BaseModel,Field
from spekforge.models.schema import SchemaSpec
from spekforge.models.param import ParamSpec

class HTTPMethod(str, Enum):
    GET= "get"
    POST= "post"
    PUT= "put"
    DELETE= "delete"
    HEAD= "head"
    OPTIONS= "options"
    PATCH= "patch"

class AuthType(str, Enum):
    NONE= "none"
    API_KEY="api_key"
    BEARER="bearer"
    BASIC="basic"
    OAUTH2="oauth2" 

class EndpointSpec(BaseModel):
    method: HTTPMethod
    path: str
    summary: str | None = None
    description: str | None = None
    parameters: list[ParamSpec] = Field(default_factory=list)
    request_body: SchemaSpec | None = None
    responses: dict[str, SchemaSpec] = Field(default_factory=dict)
    auth: AuthType = AuthType.NONE
    errors: list[str] = Field(default_factory=list)
\end{lstlisting}

\texttt{EndpointSpec} is the central component of the intermediate model, so
Phase 2 is responsible for populating the model for each endpoint detected in the documentation,
and Phase 3 uses this to build the final model with the specified
configuration. The \texttt{responses} field is modeled directly as an
extension of \texttt{dict[str, SchemaSpec]}, with the HTTP status code
as the key, mapping directly to the response block in the specification.
Both \texttt{request\_body} and \texttt{auth} are optional modules that are populated
since not all endpoints have explicit authentication or a request body.

\section{Deliverable Validation}

\subsection{Tests}

\begin{lstlisting}[language=Python, caption={\texttt{tests/unit/test\_models.py}}]
from spekforge.models.endpoint import EndpointSpec, HTTPMethod, AuthType
from spekforge.models.param import ParamSpec, ParamLocation
from spekforge.models.schema import SchemaSpec

def test_endpoint_spec_minimal():
    endpoint = EndpointSpec(method=HTTPMethod.GET, path="/pets/{id}")

    assert endpoint.method == HTTPMethod.GET
    assert endpoint.path == "/pets/{id}"
    assert endpoint.auth == AuthType.NONE
    assert endpoint.parameters == []

def test_endpoint_spec_with_params_and_responses():
    pet_schema = SchemaSpec(
        name="Pet",
        properties={
            "id": SchemaSpec(type="integer"),
            "name": SchemaSpec(type="string"),
        },
        required=["id", "name"],
    )
    id_param= ParamSpec(name="id", location=ParamLocation.PATH, required=True, type="integer")

    endpoint = EndpointSpec(
        method=HTTPMethod.GET,
        path="/pets/{id}",
        parameters=[id_param],
        responses={"200": pet_schema},
    )

    assert endpoint.parameters[0].name == "id"
    assert endpoint.responses["200"].properties["name"].type == "string"
\end{lstlisting}

The first test covers a minimum case required by the phase deliverable.
This test verifies the functionality of the intermediate model; in addition,
a second test is performed to verify the nesting of \texttt{ParamSpec} and
\texttt{SchemaSpec} within \texttt{EndpointSpec}, including the recursion
of \texttt{SchemaSpec.properties}.

\subsection{Execution}

\realcapture{0.85}{Tests.png}{\texttt{pytest} output showing the two passing tests of the intermediate model}

\realcapture{0.85}{Structure.png}{Repository file tree on GitHub, showing the complete structure of \texttt{src/spekforge/} at the close of Phase 0}


